\section{The almost Mathieu upper law down to \texorpdfstring{$L/2$}{L/2}}\label{sec:libration}

The almost Mathieu operator admits a sharper upper bound than
Corollary~\ref{cor:typeI} below $L$. Throughout this section,
$L=\log\lambda>0$. No parity assumption is needed for the upper bound.

\begin{theorem}[The sharper almost Mathieu upper law]\label{thm:libration}
Let $\alpha$ be irrational and $L/2<\beta(\alpha)<L$. For every
$a>1-L/(2\beta(\alpha))$ and $m>0$,
\[
 \sup_{\theta\in\T}P_\theta(T^a,T)=O(T^{-m}).
\]
The constants can depend on $\lambda,\alpha,a,m$.
Consequently $\aunif\le1-L/(2\beta(\alpha))$.
\end{theorem}

Together with Corollary~\ref{cor:typeI}, this proves the upper bound
$1-L/(2\beta)$ for every finite $\beta>L/2$. The matching lower bound
in Section~\ref{sec:sharp} uses $\beta_e$, so equality in terms of
$\beta$ still requires $\beta_e=\beta$.

\subsection{Reduction to a quantitative frame package}

At one convergent write
\begin{equation}\label{eq:libscale}
 q=q_n,\quad p=p_n,\quad \delta=q\alpha-p,\quad
 b=-q^{-1}\log|\delta|,\quad \Lambda=\lambda^q,\quad
 R_*=2\Lambda+2,\quad \kappa=(2\pi q)^2\Lambda\delta^2.
\end{equation}
The letter $P$ below denotes the Chambers polynomial $P_q$.
A cell has $q$ physical sites. Fix $\xi>0$, put $\tau=\xi/100$,
and choose $h=L+d_h>L$ after the finite list of derivative estimates
used below. A factor denoted $\mathcal R_q$ has the form
$q^C e^{C d_hq}$. Different occurrences can have different fixed $C$.
We first bound all these constants, then choose $d_h$ so that every
such loss fits its assigned fraction of $\tau q$, and finally take
$q$ large. Neither $d_hq$ nor $\tau q$ is called $o(q)$ at fixed
positive $d_h,\tau$.

The frame phase interval has half-width
$t_{\rm fr}=e^{-\xi q/4}/(2\pi q)$.
Its translated copies have physical centers and local phases as in
\cite[Appendix C, Covariance and well density]{DMband}.
The upper-law argument uses their raw frame estimates at
$2\xi\le b\le L-\xi$. This range includes $b\le L/2$.
It does not use a small-$\kappa$ eigenvalue-counting theorem there.

\begin{definition}[The package $\mathrm{FP}^*$]\label{def:FPstar}
For each inner edge class and scale, the following data are required,
uniformly in the phase, the well and the launch site.
\begin{enumerate}
\item A finite positive energy unit $\epsilon$, a covariant
orthonormal frame $V$, its projection $\Pi$, and summable physical
kernel bounds for the frame and its orthogonalization. At a closed
pair the unit must be defined without dividing by $P'(E_a)$.
\item A gap $m$ for $QHQ-E$ on $Q=1-\Pi$ throughout the class window.
After declaring $A=D\epsilon$, the bulk compression error $\gamma$
and bulk coupling $c_0$ satisfy
\[
 \gamma\le A/100,\qquad c_0^2/m\le A/100,\qquad
 c_0=o(m),\qquad A=o(m).
\]
A coupling at the ends of a frame interval may instead be kept as a
separate localized block. It must satisfy the boundary calculation
of Appendix~\ref{app:libboundary}.
\item In each fine energy window, at most a fixed number of shell
intervals per well and a bounded self-adjoint frame penalty give a
frame gap at least $3A$. Their widths and forbidden paddings in cells
are at most
\begin{equation}\label{eq:FPwidth}
 C\exp(C\xi q)\bigl(1+\kappa^{-1/2}\bigr).
\end{equation}
The constant $C$ in the exponent is fixed before $\xi$ is chosen.
\item For the physical operator with this penalty, a capped
cell-constant weight $\varphi$ gives a perturbation $\Delta$ with
\[
 \|\Pi\Delta\Pi\|=o(A),\quad
 \|Q\Delta\Pi\|^2+\|\Pi\Delta Q\|^2=o(mA),\quad
 \|Q\Delta Q\|=o(m).
\]
Its jumps are at the specified bonds. With $p_{\rm pad}$ cells of
forbidden padding, $\mu p_{\rm pad}\ge e^{c\xi q}$ for a fixed $c>0$.
The same bounds hold in the opposite weight direction.
\item Smooth fine windows cover the full class and overlap the
interior thresholds below. There are $e^{O(q)}$ fine windows and
$O(q)$ coarse classes. All fixed-order filter derivatives, inverse
window widths and inverse constants are $e^{O(q)}$.
\end{enumerate}
The little-$o$ comparisons refer to $q\to\infty$ after the slack
parameters have been fixed.
\end{definition}

This is a static package. In particular, the shell gap is a statement
about the linear energy-space operator $H-E$. A scalar identity for
$P(H)$ does not replace it. The factors $D$ are allowed to grow at low
$b$. Their growth must still imply~\eqref{eq:FPwidth}.

\begin{proposition}[From the package to transport]\label{prop:FPtransport}
Suppose $\mathrm{FP}^*$ holds at every inner edge class for all
sufficiently large convergents with $2\xi\le b\le L-\xi$, for every
sufficiently small $\xi>0$. Then Theorem~\ref{thm:libration} holds.
\end{proposition}

\begin{proof}
We give the scale and energy bookkeeping. The resolvent, functional
calculus and launch-shell argument are proved in
Appendix~\ref{app:liblocal}. They turn the package into absolute
filtered-packet tails smaller than $e^{-Mq}$ for every fixed $M$,
on every fixed time window $|t|\le e^{Bq}$. Their physical radius is
\begin{equation}\label{eq:libedgeR}
 R_{\rm edge}\le Cq e^{C\xi q}(1+\kappa^{-1/2}),\qquad
 q\kappa^{-1/2}=(2\pi)^{-1}e^{(b-L/2)q}.
\end{equation}
The top edge has the one-band package from the translated-well
construction in \cite[Appendices C--E]{DMband}, with the
cell-weight and penalty estimates in
Appendix~\ref{app:liblocal}. The bottom edge follows from
$u(j)\mapsto(-1)^ju(j)$ and $\theta\mapsto\theta+1/2$.

For interior energies set
$e_P(E)=\min\{|P(E)-2\Lambda|,|P(E)+2\Lambda|\}$.
If $0<\eta\le e_P/4$, then
\begin{equation}\label{eq:libsublevel}
 \bigl|\{x\in\T:|P(E)-2\Lambda\cos2\pi x|<\eta\}\bigr|
 \le \frac{2\eta}{\pi\sqrt{\Lambda e_P}}.
\end{equation}
The set has at most two arcs, up to the circle seam; it is empty if
$|P(E)|>2\Lambda$. On each monotone cosine branch its derivative is
at least $2\pi\sqrt{\Lambda e_P}$, which proves the assertion.
Replace only the sublevel step in Proposition~\ref{prop:Comega} by
this estimate. Retaining the rational-product residual gives
\begin{equation}\label{eq:libinterior}
 R_{\rm int}(E)\le
 C e^{(L/2+4\xi_C)q}e_P(E)^{-1/2}+Cq e^{\xi_Cq},
 \qquad e_P\ge16e^{(L-b+4\xi_C)q}.
\end{equation}
Here $\xi_C>0$ includes the Lyapunov-growth slack in the shrinking
spectral neighborhood. The resolvent tail is the absolute bound
of Proposition~\ref{prop:Comega}. It holds for any fixed exponential
time window. Appendix~\ref{app:liblocal} passes it to smooth interior
filters on the same type of window. No doubly exponential time
claim is needed for the interior packet.

Fix a requested exponent $a>1-L/(2\beta)$. Choose
\[
 1-L/(2\beta)<a_0<\min\{a,1/2\},\qquad
 \nu=L/2,\quad t_v=\nu/(1-a_0),\quad r_v=a_0t_v=t_v-\nu.
\]
Then $\beta<t_v<L$ and $r_v>0$. Choose
$\beta<B_0<t_v$. In the covering of Theorem~\ref{thm:U}, choose
$K$ first, and then $\xi,\xi_C,d_h$ small enough that all radius
losses fit inside the strict gaps $t_v-B_0$ and $a-a_0$.
We may require $2\xi<r_v$ and $B_0<L-\xi$.
Case A of that proof is unchanged.

In Case B put $t=\log T/q$. The three regimes are
\[
 \begin{array}{c|l}
  t\le t_v & \text{velocity},\\
  t\ge t_v,\ b\le2r_v & \text{old confinement},\\
  t\ge t_v,\ b>2r_v & \text{frames and interior}.
 \end{array}
\]
The velocity radius is
$\exp(((t-\nu)_++\text{chosen loss})q)$, up to polynomial factors.
For $t\le t_v$ it is bounded by $T^{a_0+\text{chosen loss}}$.
For $t\ge t_v$, a radius exponent at most $r_v$ suffices.
If $b\le2r_v$ the old radius exponent $b/2$ in
Proposition~\ref{prop:Camo} suffices, with the final slack assigned
to $a-a_0$. Otherwise $b>2r_v>2\xi$ and $b\le B_0<L-\xi$.
Apply the frame package. The exponent in~\eqref{eq:libedgeR} is at
most $(B_0-\nu)_++C\xi<r_v$.
For the interior use the threshold
\[
 e_P\ge \exp((2L-2t_v+8\xi_C)q).
\]
Since $b>2r_v=2t_v-L$, this dominates the threshold in
\eqref{eq:libinterior} after arbitrarily small slack. That radius
is at most $e^{(r_v+\text{chosen loss})q}$.

The edge windows extend to $P$-depth $4\Lambda e^{-\xi q}$.
Choose their plateaux strictly between the last interior threshold
and this cap. This is possible because $2L-2t_v<L$.
On $b\ge L/2+\xi$ one may take plateau
$\Lambda e^{-5\xi q/4}$, with $\xi_C\le\xi/16$.
If elementary cutoffs are $\chi_1,\ldots,\chi_k$, use
\[
 g_j=\chi_j\prod_{i<j}(1-\chi_i),\qquad
 g_{\rm rest}=\prod_{i=1}^k(1-\chi_i).
\]
These sum to one. The remainder is supported in the interior class,
after inserting a smooth spectral-neighborhood cutoff. Fixed-order
derivatives and the number of fine windows cost $e^{O(q)}$.
Overlapping windows require no false separation from the resolvent
energy. Their functional-calculus factors have $e^{O(q)}$ bounds
on $T\le e^{Bq}$.

Summing absolute packet tails costs only the number of filters.
For the Abel average, truncate time at $kqT$; the omitted part is
$e^{-2kq}$ and $kqT$ is still in a fixed exponential window.
Choose $M,k$ for the prescribed inverse power of $T$.
Far energies are treated by Combes--Thomas exactly as in
Theorem~\ref{thm:U}. Since $q$ is comparable with $\log T$ in
Case B, the resulting bound is $O(T^{-m})$ at $T^a$.
All choices are uniform in phase and launch site.
\end{proof}

\begin{remark}[The lower-ratio scales]\label{rem:liblow}
At the limiting target $a_0=1-L/(2\beta)$, the range requiring new
confinement starts at $b>2\beta-L$. For $\beta<3L/4$ it includes
$b<L/2$. This is why the raw frame construction and its coarse
energy unit are needed down to $b\ge2\xi$.
\end{remark}
